\documentclass[11pt,letterpaper]{article}
\usepackage[utf8]{inputenc}
\usepackage[T1]{fontenc}
\usepackage[spanish]{babel}
\usepackage{geometry}
\geometry{margin=2.5cm}
\usepackage{enumitem}
\setlength{\parindent}{0pt}
\setlength{\parskip}{0.7em}
\pagestyle{empty}

% Campos por completar antes de firmar: los que están entre corchetes.

\begin{document}

Bogotá D.C., [día] de noviembre de 2026

\vspace{1em}
Señores\\
\textbf{Coordinación del Programa Académico de Física}\\
\textbf{Consejo Curricular}\\
Facultad de Ciencias Matemáticas y Naturales\\
Universidad Distrital Francisco José de Caldas\\
Ciudad

\vspace{0.5em}
\textbf{Asunto:} Aval y radicación de la propuesta de Trabajo de Grado del estudiante Juan Sebastián Sánchez Álvarez.\\
\textbf{Referencia:} Comunicación PAF\_2463 del 19 de noviembre de 2026.

Cordial saludo.

En mi calidad de docente director, y de conformidad con el Artículo~4, literal~b, y el Artículo~10, literales~a y~c, del
Acuerdo N.\textsuperscript{o} 004 del 18 de agosto de 2026 del Consejo Académico, me permito avalar y radicar ante el
Consejo Curricular la siguiente propuesta de Trabajo de Grado:

\begin{itemize}[leftmargin=1.5em, itemsep=0.2em]
  \item \textbf{Título:} Análisis del problema inverso de la ecuación de Schrödinger estacionaria con potencial vectorial:
        formulación hidrodinámica y estudio de estados de Fock--Darwin.
  \item \textbf{Estudiante:} Juan Sebastián Sánchez Álvarez, código [código estudiantil].
  \item \textbf{Modalidad:} Monografía de investigación teórica (Artículo~2, literal~b, del Acuerdo 004 de 2026).
  \item \textbf{Director:} John Hernán Díaz Forero, docente [de carrera / TCO / MTO] del Programa Académico de Física.
  \item \textbf{Periodos previstos:} Trabajo de Grado en un solo periodo, 2027-1, con cierre previsto entre finales de mayo y mediados de junio de 2027.
\end{itemize}

He revisado la propuesta y considero que es pertinente para el perfil de egreso del programa, que su estructura se ajusta
a lo dispuesto en el Artículo~24 del Acuerdo para la modalidad de Monografía, y que sus objetivos, metodología y
cronograma son viables en el plazo establecido en el Artículo~9. El estudiante ha aprobado al menos el 70\,\% de los
créditos académicos del plan de estudios, según lo exige el Artículo~3. Por tratarse de un trabajo teórico y
computacional, la propuesta no involucra la participación de seres humanos ni de animales, ni el tratamiento de datos
personales.

En consecuencia, solicito comedidamente incluir la propuesta en el orden del día de la sesión del Consejo Curricular del
24 de noviembre de 2026, en la cual el estudiante, con mi acompañamiento, realizará la presentación de diez minutos
establecida por el Consejo, con el fin de inscribir el espacio académico Trabajo de Grado en el periodo 2027-1.

Se adjuntan:
\begin{enumerate}[leftmargin=1.8em, itemsep=0.1em]
  \item Documento de la propuesta de Trabajo de Grado.
  \item Presentación para la sesión del Consejo Curricular.
\end{enumerate}

Agradezco de antemano su atención.

Cordialmente,

\vspace{3.5em}
\rule{7cm}{0.4pt}\\
\textbf{John Hernán Díaz Forero}\\
Docente director\\
Programa Académico de Física\\
Facultad de Ciencias Matemáticas y Naturales\\
Universidad Distrital Francisco José de Caldas\\
Correo: [correo institucional]

\vspace{2em}
\rule{7cm}{0.4pt}\\
\textbf{Juan Sebastián Sánchez Álvarez}\\
Estudiante, código [código estudiantil]\\
Correo: [correo institucional]

\end{document}
